% GENERATED FILE. Edit canonical lessons, then run npm run generate:books.
% canonical-source-hash: fnv1a64:de48164b540a3735
% canonical-lessons: ML-C391-savasamskaram, ML-C391-smasanam, ML-C391-savappetti, ML-C391-dahippikkuka, ML-C391-kuliccituka

\chapter{A Funeral}
\label{ch:ml-a-funeral}
\hlchaptermodality{391}

\begin{chapteropening}
\textbf{By the end of this chapter:} \emph{I can say a funeral, a cremation ground, a coffin, to cremate and to bury.}

Complete the last lesson of chapter 391: I can say a funeral, a cremation ground, a coffin, to cremate and to bury.
\end{chapteropening}

\section[\texorpdfstring{\ml{ശവസംസ്കാരം} (śavasaṁskāraṁ)}{ശവസംസ്കാരം (śavasaṁskāraṁ)}]{\ml{ശവസംസ്കാരം} (śavasaṁskāraṁ) — a funeral}

\label{lesson:ML-C391-savasamskaram}



\begin{quote}
Before the new one: say the Malayalam for to say goodbye, then the Malayalam for to leave one's homeland.
\end{quote}

\subsection*{You'll want to know: \ml{ശവസംസ്കാരം}}
\textbf{\ml{ശവസംസ്കാരം}} — \emph{śavasaṁskāraṁ} — \textquotedblleft{}a funeral\textquotedblright{}.

\begin{quote}
\ml{മുത്തച്ഛന്റെ} \ml{ശവസംസ്കാരം} \ml{ഇന്നലെ} \ml{വൈകുന്നേരം} \ml{നടന്നു}.
\end{quote}

\begin{quote}
\emph{muttacchanṟe śavasaṁskāraṁ innale vaikunnēraṁ naṭannu.}
\end{quote}

\begin{quote}
Grandfather's funeral took place yesterday evening.
\end{quote}

\ml{ശവം}, a body, + \ml{സംസ്കാരം}, rites: the last rites.

\begin{cousinweb}
\textbf{\ml{ശവസംസ്കാരം}} (\emph{śavasaṁskāraṁ}) is built from two Sanskrit words, \emph{śava}, \textquotedblleft{}a corpse\textquotedblright{} (CDIAL 12356), and \emph{saṁskāra}, \textquotedblleft{}preparation, a rite\textquotedblright{} (CDIAL 13038), which also gave \ml{സംസ്കാരം} (\emph{saṁskāraṁ}), \textquotedblleft{}culture\textquotedblright{}.
\end{cousinweb}

\begin{grammarlens}[title={What you've built}]
One more word for telling what happened.
\end{grammarlens}

\subsection*{Your turn}
\begin{itemize}
\raggedright
  \item \emph{Say it:} \emph{śavasaṁskāraṁ}
  \item \emph{Say it:} \emph{śavasaṁskāraṁ}, once more
  \item \emph{Say it:} \emph{muttacchanṟe śavasaṁskāraṁ innale vaikunnēraṁ naṭannu.} — Grandfather's funeral took place yesterday evening.
  \item \emph{Recall:} say the Malayalam for to say goodbye, then the Malayalam for to leave one's homeland, then say \emph{śavasaṁskāraṁ} again
\end{itemize}

\begin{tcolorbox}[breakable,colback=teal!4,colframe=teal!35!black,title={Before you move on}]
What does \ml{ശവസംസ്കാരം} mean? (\textquotedblleft{}A funeral\textquotedblright{}.) Say it once more.
\end{tcolorbox}
\section[\texorpdfstring{\ml{ശ്മശാനം} (śmaśānaṁ)}{ശ്മശാനം (śmaśānaṁ)}]{\ml{ശ്മശാനം} (śmaśānaṁ) — a cremation ground}

\label{lesson:ML-C391-smasanam}



\begin{quote}
Before the new one: say the Malayalam for to leave one's homeland, then the Malayalam for a funeral.
\end{quote}

\subsection*{You'll want to know: \ml{ശ്മശാനം}}
\textbf{\ml{ശ്മശാനം}} — \emph{śmaśānaṁ} — \textquotedblleft{}a cremation ground\textquotedblright{}.

\begin{quote}
\ml{ഗ്രാമത്തിന്റെ} \ml{അറ്റത്ത്} \ml{ഒരു} \ml{ശ്മശാനം} \ml{ഉണ്ട്}.
\end{quote}

\begin{quote}
\emph{grāmattinṟe aṟṟattŭ oru śmaśānaṁ uṇṭŭ.}
\end{quote}

\begin{quote}
There is a cremation ground at the edge of the village.
\end{quote}

Also any burial ground.

\begin{cousinweb}
\textbf{\ml{ശ്മശാനം}} (\emph{śmaśānaṁ}) is borrowed from Sanskrit \emph{śmaśāna}, \textquotedblleft{}a place where the dead are burnt\textquotedblright{} (CDIAL 12658).
\end{cousinweb}

\begin{grammarlens}[title={What you've built}]
One more word for telling what happened.
\end{grammarlens}

\subsection*{Your turn}
\begin{itemize}
\raggedright
  \item \emph{Say it:} \emph{śmaśānaṁ}
  \item \emph{Say it:} \emph{śmaśānaṁ}, once more
  \item \emph{Say it:} \emph{grāmattinṟe aṟṟattŭ oru śmaśānaṁ uṇṭŭ.} — There is a cremation ground at the edge of the village.
  \item \emph{Recall:} say the Malayalam for to leave one's homeland, then the Malayalam for a funeral, then say \emph{śmaśānaṁ} again
\end{itemize}

\begin{tcolorbox}[breakable,colback=teal!4,colframe=teal!35!black,title={Before you move on}]
What does \ml{ശ്മശാനം} mean? (\textquotedblleft{}A cremation ground\textquotedblright{}.) Say it once more.
\end{tcolorbox}
\section[\texorpdfstring{\ml{ശവപ്പെട്ടി} (śavappeṭṭi)}{ശവപ്പെട്ടി (śavappeṭṭi)}]{\ml{ശവപ്പെട്ടി} (śavappeṭṭi) — a coffin}

\label{lesson:ML-C391-savappetti}



\begin{quote}
Before the new one: say the Malayalam for a funeral, then the Malayalam for a cremation ground.
\end{quote}

\subsection*{You'll want to know: \ml{ശവപ്പെട്ടി}}
\textbf{\ml{ശവപ്പെട്ടി}} — \emph{śavappeṭṭi} — \textquotedblleft{}a coffin\textquotedblright{}.

\begin{quote}
\ml{ആറ്} \ml{പേർ} \ml{ശവപ്പെട്ടി} \ml{പള്ളിയിലേക്ക്} \ml{ചുമന്നു}.
\end{quote}

\begin{quote}
\emph{āṟŭ pēr śavappeṭṭi paḷḷiyilēkkŭ cumannu.}
\end{quote}

\begin{quote}
Six men carried the coffin to the church.
\end{quote}

\ml{ശവം}, a body, + \ml{പെട്ടി}, a box.

\begin{cousinweb}
\textbf{\ml{ശവപ്പെട്ടി}} (\emph{śavappeṭṭi}) joins Sanskrit \emph{śava}, \textquotedblleft{}a corpse\textquotedblright{} (CDIAL 12356), and the native \ml{പെട്ടി} (\emph{peṭṭi}), \textquotedblleft{}a box\textquotedblright{} (DEDR 4388).
\end{cousinweb}

\begin{grammarlens}[title={What you've built}]
One more word for telling what happened.
\end{grammarlens}

\subsection*{Your turn}
\begin{itemize}
\raggedright
  \item \emph{Say it:} \emph{śavappeṭṭi}
  \item \emph{Say it:} \emph{śavappeṭṭi}, once more
  \item \emph{Say it:} \emph{āṟŭ pēr śavappeṭṭi paḷḷiyilēkkŭ cumannu.} — Six men carried the coffin to the church.
  \item \emph{Recall:} say the Malayalam for a funeral, then the Malayalam for a cremation ground, then say \emph{śavappeṭṭi} again
\end{itemize}

\begin{tcolorbox}[breakable,colback=teal!4,colframe=teal!35!black,title={Before you move on}]
What does \ml{ശവപ്പെട്ടി} mean? (\textquotedblleft{}A coffin\textquotedblright{}.) Say it once more.
\end{tcolorbox}
\section[\texorpdfstring{\ml{ദഹിപ്പിക്കുക} (dahippikkuka)}{ദഹിപ്പിക്കുക (dahippikkuka)}]{\ml{ദഹിപ്പിക്കുക} (dahippikkuka) — to cremate}

\label{lesson:ML-C391-dahippikkuka}



\begin{quote}
Before the new one: say the Malayalam for a cremation ground, then the Malayalam for a coffin.
\end{quote}

\subsection*{You'll want to know: \ml{ദഹിപ്പിക്കുക}}
\textbf{\ml{ദഹിപ്പിക്കുക}} — \emph{dahippikkuka} — \textquotedblleft{}to cremate\textquotedblright{}.

\begin{quote}
\ml{അവർ} \ml{മുത്തശ്ശിയുടെ} \ml{ശരീരം} \ml{വൈകുന്നേരം} \ml{ദഹിപ്പിച്ചു}.
\end{quote}

\begin{quote}
\emph{avar muttaśśiyuṭe śarīraṁ vaikunnēraṁ dahippiccu.}
\end{quote}

\begin{quote}
They cremated grandmother's body in the evening.
\end{quote}

Past: \ml{ദഹിപ്പിച്ചു}, \emph{dahippiccu}.

A causative verb built on a Sanskrit root meaning to burn.

\begin{cousinweb}
\textbf{\ml{ദഹിപ്പിക്കുക}} (\emph{dahippikkuka}) is built from Sanskrit \emph{dahana}, \textquotedblleft{}burning by fire\textquotedblright{} (CDIAL 6246), from the root \emph{dah}, \textquotedblleft{}to burn\textquotedblright{}, with the Malayalam causative ending -\ml{ിപ്പിക്കുക}.
\end{cousinweb}

\begin{grammarlens}[title={What you've built}]
One more word for telling what happened.
\end{grammarlens}

\subsection*{Your turn}
\begin{itemize}
\raggedright
  \item \emph{Say it:} \emph{dahippikkuka}
  \item \emph{Say it:} \emph{dahippikkuka}, once more
  \item \emph{Say it:} \emph{avar muttaśśiyuṭe śarīraṁ vaikunnēraṁ dahippiccu.} — They cremated grandmother's body in the evening.
  \item \emph{Recall:} say the Malayalam for a cremation ground, then the Malayalam for a coffin, then say \emph{dahippikkuka} again
\end{itemize}

\begin{tcolorbox}[breakable,colback=teal!4,colframe=teal!35!black,title={Before you move on}]
What does \ml{ദഹിപ്പിക്കുക} mean? (\textquotedblleft{}To cremate\textquotedblright{}.) Say it once more.
\end{tcolorbox}
\section[\texorpdfstring{\ml{കുഴിച്ചിടുക} (kuḻicciṭuka)}{കുഴിച്ചിടുക (kuḻicciṭuka)}]{\ml{കുഴിച്ചിടുക} (kuḻicciṭuka) — to bury}

\label{lesson:ML-C391-kuliccituka}



\begin{quote}
Before the new one: say the Malayalam for a coffin, then the Malayalam for to cremate.
\end{quote}

\subsection*{You'll want to know: \ml{കുഴിച്ചിടുക}}
\textbf{\ml{കുഴിച്ചിടുക}} — \emph{kuḻicciṭuka} — \textquotedblleft{}to bury\textquotedblright{}.

\begin{quote}
\ml{കള്ളൻ} \ml{സ്വർണ്ണം} \ml{പറമ്പിൽ} \ml{കുഴിച്ചിട്ടു}.
\end{quote}

\begin{quote}
\emph{kaḷḷan svarṇṇaṁ paṟampil kuḻicciṭṭu.}
\end{quote}

\begin{quote}
The thief buried the gold in the grounds.
\end{quote}

Past: \ml{കുഴിച്ചിട്ടു}, \emph{kuḻicciṭṭu}.

\ml{കുഴിച്ച്}, having dug, + \ml{ഇടുക}, to put: to dig and put in.

\begin{cousinweb}
\textbf{\ml{കുഴിച്ചിടുക}} (\emph{kuḻicciṭuka}) is a native Dravidian compound of \ml{കുഴിക്കുക}, \textquotedblleft{}to dig\textquotedblright{}, from \ml{കുഴി}, \textquotedblleft{}a pit, a hole\textquotedblright{} (DEDR 1818), and \ml{ഇടുക}, \textquotedblleft{}to put, to place\textquotedblright{} (DEDR 442).
\end{cousinweb}

\begin{grammarlens}[title={What you've built}]
One more word for telling what happened.
\end{grammarlens}

\subsection*{Your turn}
\begin{itemize}
\raggedright
  \item \emph{Say it:} \emph{kuḻicciṭuka}
  \item \emph{Say it:} \emph{kuḻicciṭuka}, once more
  \item \emph{Say it:} \emph{kaḷḷan svarṇṇaṁ paṟampil kuḻicciṭṭu.} — The thief buried the gold in the grounds.
  \item \emph{Recall:} say the Malayalam for a coffin, then the Malayalam for to cremate, then say \emph{kuḻicciṭuka} again
\end{itemize}

\begin{tcolorbox}[breakable,colback=teal!4,colframe=teal!35!black,title={Before you move on}]
What does \ml{കുഴിച്ചിടുക} mean? (\textquotedblleft{}To bury\textquotedblright{}.) Say it once more.
\end{tcolorbox}
